\scanpage{98}
\uline{Proposition} Si $b=c=b'=c'=0$, alors $M[b,\deleted{\unsure{b}}c,b',c';t]=M(0,0,0,0,t)$ est Azumaya, sss \deleted{$t\in k^*$.} $Q(0,0,0,0,t)=t^2\in\deleted{\unsure{b}}k^\times$ i.e. $t\in k^\times$

En effet, si $t\in k^*$, alors on ``réalise'' $t$ dans l'algèbre de matrices $M_2$ p. ex. par $u=\begin{pmatrix}0&1\\0&0\end{pmatrix}$ et $v^{\unsure{\alpha}}=t\begin{pmatrix}0&0\\1&0\end{pmatrix}$ on trouve $\operatorname{Tr}uv=t$ OK.

Si $t$ \unsure{non} inv., \unsure{la|le} \deleted{\unsure{qui}} précède montre que ça ne se réalise pas dans Azumaya.

% slanted
\margin{Il faudrait ici montrer \inserted{(\unsure{déduire du cas particulier|éclaircir ce cas particulier})} que ce critère, en termes de $Q(\text{---})\in k^\times$, est valable sans supposer $u^2=v^2=0$ (avec la seule condition \deleted{\unsure{si} $\delta\delta'\in k^{\unsure{?}}$} \deleted{$\delta=0,\delta'=0$} $\delta\delta'=0$ -- de ttes façons cela se réduit au cas d'un corps par \unsure{changt} de base}

Reste le cas où $\deleted{\unsure{b}}\delta^{\unsure{2|'}}\ne0\quad\delta'=0$

\deleted{\unsure{Ici}} On peut supposer [un peu court !!]
\[
b'=c'=0
\]
donc voici encore
\[
Q(b,c,b',c',t)=t^2
\]
donc $Q$ inv. signifie $t$ inv.

Montrons que cela signifie que $M(b,c,b',c';t)$ est Azumaya.
% archive sheet number: 97
